In 2019 an international collaboration led by Polish astronomers measured,
with very high precision and accuracy, the distance to the Large Magellanic
Cloud (LMC), a satellite galaxy of the Milky Way. In this way they set the
zero point of the extragalactic distance scale, which allowed for a very
precise measurement of the Hubble constant. Their method involved measuring
the distances to 20 eclipsing binary stars in the LMC, using the concept of
the surface brightness $S_V$ of a star defined as:
\[ S_V = m_V + 5 \log_{10} \theta, \]
where $m_V$ is the magnitude of a star in the optical $V$ band and $\theta$ is
the angular diameter of the star on the sky in milliarcseconds (mas).

The quantity $S_V$ can be understood as the magnitude of a star with an
angular diameter of 1\,mas. An empirical relation has been established between
$S_V$ and the colour index $(m_V - m_K)$, where $m_V$ and $m_K$ are magnitudes
in the $V$-band and infrared $K$-band. This is shown in the figure below for
giant stars of spectral types G and K.

\ioaafig{2023_DA_Q1-f1.pdf}

Using this relation, the distance to an eclipsing binary system can be
determined by deriving the physical radii of the components (using photometry
and spectroscopy), and comparing these with the angular diameters predicted by
the $S_V - (m_V - m_K)$ relation.

The table below gives the parameters of three eclipsing binary stars. $R_1$
and $R_2$ are the radii of each component, $V_{1+2}$ and $K_{1+2}$ are the
total brightness in magnitudes of the binary in the $V$- and $K$-bands, and
$L_2/L_1$ is the luminosity ratio of the components in each band.

\begin{center}
\begin{tabular}{l|c|c|c|c|c|c}
source ID & $R_1$ [$R_\odot$] & $R_2$ [$R_\odot$] & $V_{1+2}$ [mag]
  & $K_{1+2}$ [mag] & $L_2/L_1$ ($V$) & $L_2/L_1$ ($K$) \\
\hline
OGLE-LMC-ECL-03160 & 17.03 & 37.42 & 16.73 & 14.10 & 2.80  & 4.23  \\
OGLE-LMC-ECL-10567 & 24.60 & 36.64 & 16.15 & 13.83 & 1.41  & 1.99  \\
OGLE-LMC-ECL-18365 & 37.30 & 15.94 & 16.27 & 14.01 & 0.206 & 0.188 \\
\end{tabular}
\end{center}

Apply the method outlined above to the three eclipsing binary systems and
calculate the distance to the LMC in kiloparsecs. Estimate the total error of
the result. Assume that the fitting of the $S_V - (m_V - m_K)$ relation
contributes to a bias of up to 0.8\% in all measurements simultaneously.

Hint: in your calculations keep at least three significant figures and two
decimal places. Assume that interstellar extinction is negligible and that the
angular size of the LMC is small.
